\documentclass[12pt,a4paper]{article}
\usepackage[utf8]{vietnam}
\usepackage[left=3.00cm, right=2.00cm, top=2.50cm, bottom=2.50cm]{geometry}
\usepackage{amsmath}
\usepackage{amsfonts}
\usepackage{amssymb}
\usepackage{graphicx}
\usepackage{float}
\usepackage{array}
\usepackage{booktabs}
\usepackage{tabularx}
\usepackage{multirow}
\usepackage{hyperref}
\usepackage{titlesec}
\usepackage{fancyhdr}
\usepackage{tikz}
\usepackage{listings}
\usepackage{xcolor}
\usetikzlibrary{calc}

\hypersetup{
    colorlinks=true,
    linkcolor=black,
    filecolor=black,      
    urlcolor=black,
}

\pagestyle{fancy}
\fancyhf{}
\rhead{\small{\textit{Báo cáo thực tập tuần 3}}}
\lfoot{\small{Sinh viên thực hiện: Nguyễn Thanh Hà}}
\rfoot{\thepage}

\titleformat{\section}{\normalfont\fontsize{14}{17}\bfseries\selectfont\color{black}}{\thesection}{1em}{}
\titleformat{\subsection}{\normalfont\fontsize{12}{15}\bfseries\selectfont\color{black}}{\thesubsection}{1em}{}
\titleformat{\subsubsection}{\normalfont\fontsize{12}{15}\itshape\selectfont\color{black}}{\thesubsubsection}{1em}{}

\lstdefinestyle{cppstyle}{
    backgroundcolor=\color{gray!10},
    commentstyle=\color{green!60!black},
    keywordstyle=\color{blue}\bfseries,
    numberstyle=\tiny\color{gray},
    stringstyle=\color{purple},
    basicstyle=\ttfamily\footnotesize,
    breaklines=true,
    captionpos=b,
    numbers=left,
    numbersep=5pt,
    frame=single
}
\lstset{style=cppstyle}

\begin{document}

% ==================== TRANG BÌA ====================
\begin{titlepage}
\begin{tikzpicture}[remember picture, overlay]
  \draw[line width = 1.5pt] ($(current page.north west) + (3.0cm,-2.0cm)$) rectangle ($(current page.south east) + (-1.5cm,2.0cm)$);
  \draw[line width = 0.5pt] ($(current page.north west) + (3.1cm,-2.1cm)$) rectangle ($(current page.south east) + (-1.6cm,2.1cm)$);
\end{tikzpicture}

\begin{center}
    \textbf{\fontsize{14}{17}\selectfont TRƯỜNG ĐẠI HỌC CÔNG NGHỆ} \\
    \textbf{\fontsize{12}{15}\selectfont ĐẠI HỌC QUỐC GIA HÀ NỘI} \\
    
    \vspace{2.5cm}
    \textbf{\fontsize{18}{22}\selectfont BÁO CÁO THỰC TẬP TUẦN 3} \\[0.5cm]
    \textbf{\fontsize{13}{16}\selectfont ĐỀ TÀI: TÁI CẤU TRÚC C++ CORE SANG REMOTE HTTP CONTROL} \\
    \textbf{\fontsize{13}{16}\selectfont \& THÊM EMBEDDED HTTP SERVER + DATA SYNC MANAGER} \\
    
    \vspace{3.5cm}
    \fontsize{12}{15}\selectfont
    \begin{tabular}{ll}
        Sinh viên thực hiện: & \textbf{Nguyễn Thanh Hà} \\
        Mã số sinh viên: & \textbf{23021810} \\
        Lớp: & \textbf{K68E-EC1} \\
        Người hướng dẫn: & \textbf{Hồ Sỹ Minh} \\
        Giáo viên hướng dẫn: & \textbf{Nguyễn Hồng Thịnh} \\
    \end{tabular}
    
    \vfill
    \textbf{Hà Nội, 2026}
\end{center}
\end{titlepage}

\newpage
\tableofcontents
\newpage

% ==================== CHƯƠNG 1 ====================
\section{MỤC TIÊU VÀ KẾ HOẠCH CHI TIẾT TUẦN 3}
Mục tiêu chính của Tuần 3 là giải quyết triệt để vấn đề treo ứng dụng console C++ do cơ chế nhập liệu qua \texttt{cin}, tái cấu trúc ứng dụng sang State Machine điều khiển từ xa bất đồng bộ (\texttt{DeviceMode}), phát triển Embedded HTTP Server (\texttt{device\_server}) lắng nghe cổng 8080 và module Data Sync Manager (\texttt{sync\_manager}) giao tiếp với Central Backend.

\begin{table}[H]
\centering
\small
\begin{tabularx}{\textwidth}{|l|X|X|}
\hline
\textbf{Thời gian} & \textbf{Nội dung công việc thực hiện} & \textbf{Kết quả / Sản phẩm yêu cầu} \\ \hline
\textbf{Thứ Hai} & Phân tích nguyên nhân lệnh \texttt{cin} làm treo luồng camera đơn. Xóa 100\% \texttt{cin} trong \texttt{main.cpp}. Khởi tạo các biến atomic: \texttt{deviceMode}, \texttt{pendingResidentId}, \texttt{enrollProgress}, \texttt{elevatorMutex}. & Vòng lặp chính C++ rẽ nhánh bất đồng bộ giữa \texttt{RECOGNIZING} và \texttt{ENROLLING}. \\ \hline
\textbf{Thứ Ba} & Cập nhật hàm \texttt{enrollFace()} (bảo toàn 100\% chữ ký hàm), bổ sung logic cập nhật tiến độ 0..30 mẫu và nhận lệnh thoát từ xa khi Cancel. & Thu thập mẫu ảnh không còn bị treo bàn phím console. \\ \hline
\textbf{Thứ Tư} & Lập trình module Embedded HTTP Server (\texttt{device\_server.h/.cpp}) dùng \texttt{cpp-httplib} chạy trên \texttt{std::thread} riêng tại port 8080. & Hoàn thành 3 REST endpoints: \texttt{start-enrollment}, \texttt{status}, \texttt{cancel-enrollment}. \\ \hline
\textbf{Thứ Năm} & Lập trình module Data Sync Manager (\texttt{sync\_manager.h/.cpp}) với HTTP Client đồng bộ dữ liệu cư dân, gửi \texttt{access-logs} và \texttt{alerts}. & Tự động đẩy nhật ký và cảnh báo an ninh về Backend. \\ \hline
\textbf{Thứ Sáu} & Tích hợp vào \texttt{main.cpp}, cập nhật \texttt{CMakeLists.txt}, thực thi 2 bộ Unit Test Suite (\texttt{test\_device\_server.cpp} \& \texttt{test\_sync\_manager.cpp}). & Tất cả 9 test cases **PASSED 100\%**, chịu lỗi ngắt mạng an toàn. \\ \hline
\end{tabularx}
\caption{Phân bổ lịch trình công việc Tuần 3}
\end{table}

\newpage

% ==================== CHƯƠNG 2 ====================
\section{NỘI DUNG VÀ KẾT QUẢ THỰC HIỆN CHI TIẾT}

\subsection{Thứ Hai — Loại bỏ cơ chế cin \& Tái cấu trúc State Machine}
\begin{itemize}
    \item \textbf{Nội dung công việc}: Phân tích luồng thực thi trong \texttt{main.cpp}. Khi ứng dụng gọi \texttt{cin >> choice} hoặc \texttt{getline(cin, input)}, HĐH Windows tạm dừng (\textit{block}) toàn bộ thread chính để chờ dữ liệu từ bàn phím. Điều này dẫn đến việc luồng hiển thị OpenCV (\texttt{imshow}) và luồng xử lý camera bị đóng băng hoàn toàn.
    \item \textbf{Giải pháp kiến trúc}: Xóa bỏ hoàn toàn câu lệnh \texttt{cin}. Chuyển sang mô hình máy trạng thái thread-safe bằng các biến atomic toàn cục:
\end{itemize}

\begin{lstlisting}[language=C++]
// State Machine kieu enum class va cac bien atomic thread-safe
enum class DeviceMode { RECOGNIZING, ENROLLING };
std::atomic<DeviceMode> deviceMode{DeviceMode::RECOGNIZING};
std::atomic<int> pendingResidentId{-1};
std::atomic<int> enrollProgress{0};
std::mutex elevatorMutex;
\end{lstlisting}

\begin{itemize}
    \item \textbf{Luồng rẽ nhánh}: Vòng lặp chính \texttt{while(true)} liên tục kiểm tra biến \texttt{deviceMode.load()}:
    \begin{enumerate}
        \item Nếu là \texttt{RECOGNIZING}: Tiến hành lấy frame camera, nhận diện khuôn mặt AI LBPH, cập nhật trạng thái thang máy.
        \item Nếu là \texttt{ENROLLING}: Lấy thông tin cư dân theo \texttt{pendingResidentId.load()}, tự động gọi hàm \texttt{enrollFace()}, thu thập 30 ảnh mẫu, sau đó huấn luyện lại mô hình LBPH và tự động chuyển về \texttt{RECOGNIZING}.
    \end{enumerate}
\end{itemize}

\subsection{Thứ Ba — Cập nhật hàm enrollFace() hỗ trợ Cancel \& Progress Tracking}
\begin{itemize}
    \item \textbf{Nội dung công việc}: Chỉnh sửa hàm \texttt{enrollFace()} để thông báo tiến độ thu thập ảnh mẫu ra bên ngoài mà vẫn **bảo toàn 100\% chữ ký hàm** như yêu cầu ban đầu:
\end{itemize}

\begin{lstlisting}[language=C++]
void enrollFace(VideoCapture& cap, CascadeClassifier& faceCascade, 
                DatabaseManager& db, const Resident& resident) {
    int sampleCount = db.getFaceImageCount(resident.id);
    enrollProgress = sampleCount;
    
    while (true) {
        // Kiem tra lenh huy tu xa
        if (deviceMode.load() != DeviceMode::ENROLLING) {
            cout << "[DANG KY] Nhan lenh huy dang ky tu xa (Cancel). Thoat." << endl;
            break;
        }
        cap >> frame;
        if (frame.empty()) break;
        
        // Neu phat hien va luu anh mau thanh cong:
        if (db.saveFaceImage(resident.id, resized, sampleCount)) {
            sampleCount++;
            enrollProgress = sampleCount; // Cap nhat tien do atomic 0..30
        }
        if (sampleCount >= 30) break;
    }
}
\end{lstlisting}

\subsection{Thứ Tư — Lập trình Module Embedded HTTP Server (device\_server)}
\begin{itemize}
    \item \textbf{Nội dung công việc}: Xây dựng lớp \texttt{DeviceServer} trong \texttt{device\_server.h} và \texttt{device\_server.cpp} dựa trên thư viện lightweight \texttt{cpp-httplib}. HTTP Server khởi chạy trên 1 \texttt{std::thread} riêng biệt tại cổng \texttt{8080}.
    \item \textbf{Chi tiết endpoints}:
\end{itemize}

\begin{lstlisting}[language=C++]
// Endpoints cua DeviceServer (cổng 8080)
1. POST /device/start-enrollment
   Body: {"residentId": 5}
   Logic: Neu deviceMode == RECOGNIZING -> Ghi pendingResidentId=5, deviceMode=ENROLLING, tra 200 OK.
          Neu deviceMode == ENROLLING   -> Tra HTTP 409 Conflict ("Device busy").

2. GET /device/status
   Response: {"mode": "ENROLLING", "enrollProgress": 15, "elevatorState": "MOVING", "currentFloor": 4, "targetFloor": 8}

3. POST /device/cancel-enrollment
   Logic: Dat deviceMode = RECOGNIZING, reset pendingResidentId = -1, tra 200 OK.
\end{lstlisting}

\subsection{Thứ Năm — Lập trình Module Data Sync Manager (sync\_manager)}
\begin{itemize}
    \item \textbf{Nội dung công việc}: Tạo lớp \texttt{SyncManager} trong \texttt{sync\_manager.h/.cpp} làm nhiệm vụ HTTP Client kết nối tới Backend trung tâm (port 3000):
    \begin{enumerate}
        \item \texttt{syncResidentsFromBackend(db)}: Gửi request \texttt{GET /api/device/residents} để tải toàn bộ danh sách cư dân mới nhất về lưu vào SQLite local.
        \item \texttt{sendAccessLog(residentId, floor)}: Gửi \texttt{POST /api/device/access-logs} ghi nhận lượt quẹt mặt thành công.
        \item \texttt{sendAlert(reason)}: Gửi \texttt{POST /api/device/alerts} thông báo khi phát hiện Người lạ.
    \end{enumerate}
    \item \textbf{Cơ chế bảo mật \& chịu lỗi}: Mỗi request đều gửi kèm HTTP Header \texttt{X-API-Key: DEV\_SECRET\_KEY\_123}. Tất cả các lệnh HTTP Client đều cài đặt \texttt{set\_connection\_timeout(1, 0)} để khi ngắt mạng, ứng dụng C++ chỉ ghi log cảnh báo và tiếp tục chạy mượt mà, **không bị crash hay treo camera**.
\end{itemize}

\subsection{Thứ Sáu — Kiểm thử Tự động \& Biên dịch CMake}
\begin{itemize}
    \item \textbf{Nội dung công việc}: Tích hợp \texttt{DeviceServer} và \texttt{SyncManager} vào \texttt{main.cpp}. Cập nhật file \texttt{CMakeLists.txt} liên kết các thư viện mạng Win32 \texttt{ws2\_32} và \texttt{crypt32}.
    \item \textbf{Biên dịch & Test}: Viết hai bộ test tự động \texttt{test\_device\_server.cpp} (kiểm tra 5 kịch bản HTTP REST) và \texttt{test\_sync\_manager.cpp} (kiểm tra 4 kịch bản sync dữ liệu với Mock Server). Biên dịch thành công \texttt{SmartElevatorCamera.exe}.
\end{itemize}

\newpage

% ==================== CHƯƠNG 3 ====================
\section{KẾT QUẢ KIỂM THỬ ĐƠN VỊ (UNIT TEST)}

\begin{table}[H]
\centering
\small
\begin{tabularx}{\textwidth}{|l|X|X|c|}
\hline
\textbf{Mã TC} & \textbf{Mục tiêu kiểm thử} & \textbf{Dữ liệu đầu vào \& Kịch bản} & \textbf{Trạng thái} \\ \hline
TC01 & Polling GET /device/status & Gọi API khi thang IDLE & Pass (200 OK) \\ \hline
TC02 & POST /device/start-enrollment & Gửi residentId = 1 khi mode RECOGNIZING & Pass (200 OK) \\ \hline
TC03 & Start Enrollment khi bận & Gửi residentId = 2 khi mode đang ENROLLING & Pass (409 Conflict) \\ \hline
TC04 & POST /device/cancel-enrollment & Gọi API khi mode đang ENROLLING & Pass (200 OK) \\ \hline
TC05 & Start Enrollment ID không có & Gửi residentId = 9999 không tồn tại & Pass (404 Not Found) \\ \hline
TC06 & Đồng bộ cư dân từ Backend & Gọi \texttt{syncResidentsFromBackend()} & Pass (Updated DB) \\ \hline
TC07 & Gửi Access Log & Gọi \texttt{sendAccessLog(1, 5)} & Pass (Log Sent) \\ \hline
TC08 & Gửi Security Alert & Gọi \texttt{sendAlert("Phat hien nguoi la")} & Pass (Alert Sent) \\ \hline
TC09 & Thử nghiệm Backend Offline & Tắt Backend Server và gọi các hàm sync & Pass (No Crash) \\ \hline
\end{tabularx}
\caption{Bảng tổng hợp kết quả kiểm thử đơn vị Tuần 3}
\end{table}

\section{MINH HỌA KẾT QUẢ HỆ THỐNG THỰC TẾ}

\begin{figure}[H]
    \centering
    \includegraphics[width=0.55\textwidth]{images/elevator_ui.png}
    \caption{Giao diện Thang Máy C++ (Smart Elevator Controller) — Hiển thị số tầng LED xanh (tầng 01), thanh trượt vị trí thang và cửa đóng mở theo thời gian thực.}
    \label{fig:elevator_ui}
\end{figure}

\begin{figure}[H]
    \centering
    \includegraphics[width=0.75\textwidth]{images/camera_ai.png}
    \caption{Hệ thống Camera AI Thang Máy (OpenCV LBPH) — Phát hiện và khoanh vùng khuôn mặt bằng khung đỏ, hiển thị nhãn nhận diện \texttt{Nguoi la! Khong co quyen [115]} khi phát hiện người lạ không có quyền truy cập.}
    \label{fig:camera_ai}
\end{figure}

\section{ĐÁNG GIÁ TIẾN ĐỘ VÀ KẾ HOẠCH TUẦN TIẾP THEO}
\begin{itemize}
    \item \textbf{Đánh giá cá nhân}: Hoàn thành 100\% mục tiêu đặt ra cho Tuần 3. Loại bỏ hoàn toàn lỗi treo console \texttt{cin}, ứng dụng C++ sẵn sàng điều khiển từ xa qua REST APIs.
    \item \textbf{Vấn đề phát sinh \& cách xử lý}: Thư viện \texttt{cpp-httplib} cần liên kết các thư viện mạng bổ sung \texttt{ws2\_32} và \texttt{crypt32} của Windows. Đã cập nhật \texttt{CMakeLists.txt} để tự động liên kết khi phát hiện môi trường Win32.
    \item \textbf{Kế hoạch Tuần 4}: Thiết kế sơ đồ CSDL 9 bảng và phát triển Central Backend Server (NestJS / Express TypeScript) đóng vai trò trung tâm điều hành toàn hệ thống.
\end{itemize}

\vspace{1.5cm}
\begin{table}[H]
    \centering
    \begin{tabular}{p{7cm} p{7cm}}
        \centering \textbf{XÁC NHẬN CỦA MENTOR} & \centering \textbf{SINH VIÊN THỰC TẬP} \\[3cm]
        \centering \textbf{(Ký và ghi rõ họ tên)} & \centering \textbf{Nguyễn Thanh Hà} \\
    \end{tabular}
\end{table}

\end{document}
